\begin{theorem}[window-\texorpdfstring{$6$}{6} 局部支撑理想的四重秩一致性]\label{thm:conclusion-window6-local-support-ideal-four-rank-consistency}
\leanverified{paper\_conclusion\_window6\_local\_support\_ideal\_four\_rank\_consistency}
对任意子集 \(Y\subseteq X_6\)，定义
\[
A_Y:=\bigoplus_{w\in Y}M_{d(w)}(\CC),
\qquad
G_Y:=\prod_{w\in Y}\mathfrak S_{d(w)}.
\]
则恒有严格等式
\[
\operatorname{rank}K_0(A_Y)
=
\dim_{\CC}Z(A_Y)
=
\operatorname{rank}_{\FF_2}\bigl(G_Y^{\mathrm{ab}}\bigr)
=
|Y|.
\]
故在 window-\(6\) 上，任意稳定类型子系的四类“尺寸”——\(K\)-理论秩、中心维数、阿贝尔电荷秩与支撑基数——完全一致。这一一致性不是全局偶然，而是逐子集成立。
\end{theorem}

\begin{proof}
由定理 \ref{thm:conclusion-window6-groupoid-collision-dimension-identity} 的 Wedderburn 分解，\(A_Y\) 正是从
\[
A_6\cong M_2(\CC)^{\oplus 8}\oplus M_3(\CC)^{\oplus 4}\oplus M_4(\CC)^{\oplus 9}
\]
中抽取 \(Y\) 所对应的简单块直和。故
\[
\dim_{\CC}Z(A_Y)=|Y|.
\]
又因为每个矩阵块都满足
\[
K_0(M_n(\CC))\cong \ZZ,
\]
所以
\[
K_0(A_Y)\cong \ZZ^{|Y|},
\qquad
\operatorname{rank}K_0(A_Y)=|Y|.
\]
另一方面，window-\(6\) 的所有纤维尺寸都属于 \(\{2,3,4\}\)，因而
\[
\mathfrak S_{d(w)}^{\mathrm{ab}}\cong \ZZ/2\ZZ
\qquad (w\in X_6).
\]
于是
\[
G_Y^{\mathrm{ab}}\cong (\ZZ/2\ZZ)^{|Y|},
\]
其 \(\FF_2\)-秩即为 \(|Y|\)。三式合并即得结论。证毕。
\end{proof}

\begin{proposition}[window-\texorpdfstring{$6$}{6} 阈值滤层的半单闭式与秩梯]\label{prop:conclusion-window6-threshold-filtration-semisimple-rank-ladder}
\leanverified{paper\_conclusion\_window6\_threshold\_filtration\_semisimple\_rank\_ladder}
对 \(r\in\{2,3,4\}\) 定义
\[
Y_{\ge r}:=\{w\in X_6:\ d(w)\ge r\},
\qquad
A_{\ge r}:=A_{Y_{\ge r}},
\qquad
G_{\ge r}:=G_{Y_{\ge r}}.
\]
则有精确分解
\[
A_{\ge2}\cong M_2(\CC)^{\oplus 8}\oplus M_3(\CC)^{\oplus 4}\oplus M_4(\CC)^{\oplus 9},
\]
\[
A_{\ge3}\cong M_3(\CC)^{\oplus 4}\oplus M_4(\CC)^{\oplus 9},
\]
\[
A_{\ge4}\cong M_4(\CC)^{\oplus 9},
\]
并且
\[
\operatorname{rank}K_0(A_{\ge2})
=
\dim Z(A_{\ge2})
=
\operatorname{rank}_{\FF_2}(G_{\ge2}^{\mathrm{ab}})
=21,
\]
\[
\operatorname{rank}K_0(A_{\ge3})
=
\dim Z(A_{\ge3})
=
\operatorname{rank}_{\FF_2}(G_{\ge3}^{\mathrm{ab}})
=13,
\]
\[
\operatorname{rank}K_0(A_{\ge4})
=
\dim Z(A_{\ge4})
=
\operatorname{rank}_{\FF_2}(G_{\ge4}^{\mathrm{ab}})
=9.
\]
因此，window-\(6\) 的尾部多重度计数 \((21,13,9)\) 在半单代数、\(K\)-理论与阿贝尔电荷三个层面同步再现。
\end{proposition}

\begin{proof}
由 window-\(6\) 的纤维直方图
\[
\#\{d=2\}=8,\qquad
\#\{d=3\}=4,\qquad
\#\{d=4\}=9,
\]
可得
\[
|Y_{\ge2}|=8+4+9=21,
\qquad
|Y_{\ge3}|=4+9=13,
\qquad
|Y_{\ge4}|=9.
\]
将这些计数直接代入 \(A_Y\) 与 \(G_Y\) 的定义，即得三条 Wedderburn 闭式。随后由定理 \ref{thm:conclusion-window6-local-support-ideal-four-rank-consistency} 应用于 \(Y_{\ge r}\) 即得对应的秩等式。证毕。
\end{proof}

\begin{theorem}[尾计数测度的离散 Abel--Stieltjes 作用公式]\label{thm:conclusion-window6-tail-count-discrete-abel-stieltjes-action}
设有限纤维尺寸直方图为 \((n_d)_{2\le d\le D}\)，并定义尾计数
\[
N(r):=\sum_{d\ge r}n_d
\qquad (2\le r\le D).
\]
则有四条精确恒等式
\[
\sum_{d=2}^{D} n_d(d-1)=\sum_{r=2}^{D} N(r),
\]
\[
\sum_{d=2}^{D} n_d\binom d2=\sum_{r=2}^{D}(r-1)N(r),
\]
\[
\sum_{d=2}^{D} n_d\log(d!)=\sum_{r=2}^{D}N(r)\log r,
\]
\[
\sum_{d=2}^{D} n_d\, d\log d
=
\sum_{r=2}^{D}N(r)\Bigl[r\log r-(r-1)\log(r-1)\Bigr].
\]
因此，一旦尾计数测度
\[
\nu:=\sum_{r=2}^{D}N(r)\,\delta_r
\]
被确定，则线性维数、反射数、规范体积与熵缺陷这四类量全部由同一离散 Stieltjes 作用统一产生。
\leanverified{paper\_conclusion\_window6\_tail\_count\_discrete\_abel\_stieltjes\_action}
\end{theorem}

\begin{proof}
逐项写成从 \(2\) 到 \(d\) 的累积和：
\[
d-1=\sum_{r=2}^{d}1,
\qquad
\binom d2=\sum_{r=2}^{d}(r-1),
\qquad
\log(d!)=\sum_{r=2}^{d}\log r.
\]
再设 \(A_r:=r\log r\) 且 \(A_1:=0\)，则
\[
d\log d=A_d=\sum_{r=2}^{d}(A_r-A_{r-1})
=\sum_{r=2}^{d}\bigl[r\log r-(r-1)\log(r-1)\bigr].
\]
将上述四式分别乘以 \(n_d\) 并对 \(d\) 求和，交换求和次序，即得
\[
\sum_d n_d\sum_{r=2}^{d}\phi(r)
=
\sum_{r=2}^{D}\phi(r)\sum_{d\ge r}n_d
=
\sum_{r=2}^{D}\phi(r)N(r),
\]
其中 \(\phi\) 依次取 \(1\)、\(r-1\)、\(\log r\) 与 \(A_r-A_{r-1}\)。证毕。
\end{proof}

\begin{corollary}[window-\texorpdfstring{$6$}{6} 的规范--熵常数闭式]\label{cor:conclusion-window6-gauge-entropy-constants-from-tail-measure}
\leanverified{paper\_conclusion\_window6\_gauge\_entropy\_constants\_from\_tail\_measure}
对 window-\(6\)，有
\[
(n_2,n_3,n_4)=(8,4,9),
\qquad
(N_2,N_3,N_4)=(21,13,9).
\]
设
\[
\mathcal R_6:=\bigoplus_{w\in X_6}\left\{a\in \CC^{F_w}:\ \sum_{n\in F_w}a_n=0\right\},
\qquad
G_6^{\mathrm{bin}}\cong \mathfrak S_2^{8}\times \mathfrak S_3^{4}\times \mathfrak S_4^{9}.
\]
则有严格闭式
\[
\dim_{\CC}\mathcal R_6=43,
\qquad
\#\mathrm{Ref}(G_6^{\mathrm{bin}})=74,
\qquad
\operatorname{rank}_{\FF_2}(G_6^{\mathrm{bin}})^{\mathrm{ab}}=21,
\]
\[
\log|G_6^{\mathrm{bin}}|
=
39\log 2+13\log 3,
\]
\[
\kappa_6=\frac{88\log 2+12\log 3}{64},
\qquad
g_6=\frac{39\log 2+13\log 3}{64},
\]
并且
\[
\kappa_6-g_6=\frac{49\log 2-\log 3}{64}.
\]
于是，window-\(6\) 的“信息亏损--规范体积”差不是数值噪声，而是尾计数测度
\[
\nu_6=21\delta_2+13\delta_3+9\delta_4
\]
对离散步长核
\[
(r-1)\log\frac{r}{r-1}
\]
的精确作用值。
\end{corollary}

\begin{proof}
由 \(n_2=8,n_3=4,n_4=9\) 立即得到
\[
N_2=21,\qquad N_3=13,\qquad N_4=9.
\]
再把定理 \ref{thm:conclusion-window6-tail-count-discrete-abel-stieltjes-action} 应用于 \(D=4\)，即可得到
\[
\dim_{\CC}\mathcal R_6=\sum_d n_d(d-1)=\sum_{r=2}^{4}N_r=43,
\]
\[
\#\mathrm{Ref}(G_6^{\mathrm{bin}})
=
\sum_d n_d\binom d2
=
\sum_{r=2}^{4}(r-1)N_r
=21+26+27=74,
\]
\[
\log|G_6^{\mathrm{bin}}|
=
\sum_d n_d\log(d!)
=
\sum_{r=2}^{4}N_r\log r
=39\log 2+13\log 3.
\]
于是
\[
g_6=\frac{1}{64}\log|G_6^{\mathrm{bin}}|
=
\frac{39\log 2+13\log 3}{64}.
\]
同理，
\[
64\,\kappa_6
=
\sum_d n_d\,d\log d
=
\sum_{r=2}^{4}N_r\Bigl[r\log r-(r-1)\log(r-1)\Bigr]
=88\log 2+12\log 3.
\]
相减即得
\[
\kappa_6-g_6=\frac{49\log 2-\log 3}{64}.
\]
阿贝尔化秩则由推论 \ref{cor:conclusion-window6-rational-homotopy-determines-gauge-volume} 给出。证毕。
\end{proof}

\begin{corollary}[window-\texorpdfstring{$6$}{6} 的 coinvariant Poincar\'e 包由尾测度生成]\label{cor:conclusion-window6-coinvariant-poincare-package-generated-by-tail-measure}
window-\(6\) 的隐藏 coinvariant 代数 \(\mathcal C_6\) 的 Hilbert 级数可重写为
\[
\operatorname{Hilb}(\mathcal C_6;t)
=
\prod_{r=2}^{4}\left(\sum_{j=0}^{r-1}t^j\right)^{N_r}
=
(1+t)^{21}(1+t+t^2)^{13}(1+t+t^2+t^3)^{9}.
\]
特别地，
\[
t^{74}\operatorname{Hilb}(\mathcal C_6;t^{-1})
=
\operatorname{Hilb}(\mathcal C_6;t),
\qquad
\dim_{\CC}\mathcal C_6=2^{39}3^{13}=|G_6^{\mathrm{bin}}|.
\]
故 window-\(6\) 的 coinvariant 结构不是一组孤立的 Weyl 因子乘积，而是尾计数测度 \(\nu_6\) 的直接 Poincar\'e 生成：其指数 \((21,13,9)\) 正是容量右导数，也正是隐藏热核的三个基本长度层。
\end{corollary}

\begin{proof}
这只是把定理 \ref{thm:conclusion-window6-hidden-coinvariant-gorenstein-package} 中的显式 Hilbert 级数改写成尾计数指数记号
\[
(N_2,N_3,N_4)=(21,13,9).
\]
回文性、顶次数 \(74\) 与总维数 \(2^{39}3^{13}\) 亦均由该定理直接给出。证毕。
\end{proof}

\begin{theorem}[对象层种子唯一锁定无标号半单型，仅差等型块置换]\label{thm:conclusion-window6-object-seed-fixes-unlabelled-semisimple-type-up-to-block-permutation}
设 \((A_6,\alpha)\) 为 window-\(6\) 的 involutive algebra 对象。则对象层三元种子
\[
\bigl(\dim_{\CC}Z(A_6),\ \dim_{\CC}A_6^{\alpha,-},\ \dim_{\CC}A_6\bigr)
=(21,64,212)
\]
唯一决定下列全部无标号结构：
\[
\mu_6=8\delta_2+4\delta_3+9\delta_4,
\]
\[
A_6\cong M_2(\CC)^{\oplus 8}\oplus M_3(\CC)^{\oplus 4}\oplus M_4(\CC)^{\oplus 9},
\]
\[
G_6^{\mathrm{bin}}\cong \mathfrak S_2^{8}\times \mathfrak S_3^{4}\times \mathfrak S_4^{9}.
\]
并且这已经是最大可能的刚性：剩余的一切不唯一性，恰好等于同型 Wedderburn 块的置换自由度
\[
\mathfrak S_8\times \mathfrak S_4\times \mathfrak S_9.
\]
因此，window-\(6\) 的对象层种子在代数上是完全刚性的，在几何上则只留下标签级不刚性。
\end{theorem}

\begin{proof}
由定理 \ref{thm:conclusion-window6-groupoid-involution-lefschetz}，
\[
\dim_{\CC}A_6^{\alpha,-}=64=|\Omega_6|.
\]
再由定理 \ref{thm:conclusion-window6-qmoment-triple-geometry} 与定理 \ref{thm:conclusion-window6-groupoid-collision-dimension-identity}，
\[
\dim_{\CC}Z(A_6)=S_0=21,
\qquad
\dim_{\CC}A_6=S_2=212.
\]
故该三元种子实际上唯一恢复了
\[
(S_0,S_1,S_2)=(21,64,212).
\]
随后，定理 \ref{thm:conclusion-window6-first-three-moments-recover-wedderburn-type} 立即推出
\[
A_6\cong M_2(\CC)^{\oplus 8}\oplus M_3(\CC)^{\oplus 4}\oplus M_4(\CC)^{\oplus 9},
\]
从而也唯一恢复
\[
\mu_6=8\delta_2+4\delta_3+9\delta_4.
\]
再由推论 \ref{cor:conclusion-window6-rational-homotopy-determines-gauge-volume}，
\[
G_6^{\mathrm{bin}}\cong \mathfrak S_2^{8}\times \mathfrak S_3^{4}\times \mathfrak S_4^{9}.
\]
至于剩余不唯一性，只可能来自同型简单块之间的置换。由
\[
\Aut_{\CC\text{-alg}}(A_6)
\cong
\bigl(\operatorname{PGL}_2(\CC)^8\rtimes \mathfrak S_8\bigr)
\times
\bigl(\operatorname{PGL}_3(\CC)^4\rtimes \mathfrak S_4\bigr)
\times
\bigl(\operatorname{PGL}_4(\CC)^9\rtimes \mathfrak S_9\bigr)
\]
可知，这种残余恰好由
\[
\mathfrak S_8\times \mathfrak S_4\times \mathfrak S_9
\]
给出。证毕。
\end{proof}

\begin{conjecture}[rigid window 的种子--半单型--定位缺口原则]\label{conj:conclusion-rigid-window-seed-semisimple-type-localization-gap-principle}
对任意 rigid dyadic window \(m\)，设其 bin-fold 纤维群胚代数为
\[
A_m=\bigoplus_{x\in X_m}M_{d_m(x)}(\CC),
\qquad
G_m=\prod_{x\in X_m}\mathfrak S_{d_m(x)}.
\]
若对象层 involutive seed 足以恢复纤维测度 \(\mu_m\)，则应同时成立：
\begin{enumerate}
  \item \(\mu_m\) 唯一决定 \(A_m\) 的无标号 Wedderburn 型与 \(G_m\) 的无标号直积分解；
  \item \(\mu_m\) 一般不能决定等型块在 \(X_m\) 中的几何定位；
  \item 对任意阈值子集
  \[
  Y_{\ge r}:=\{x\in X_m:\ d_m(x)\ge r\},
  \]
  都有
  \[
  \operatorname{rank}K_0(A_{m,\ge r})
  =
  \dim Z(A_{m,\ge r})
  =
  \operatorname{rank}(G_{m,\ge r}^{\mathrm{ab}})
  =
  N_m(r),
  \]
  其中
  \[
  A_{m,\ge r}:=\bigoplus_{d_m(x)\ge r}M_{d_m(x)}(\CC),
  \qquad
  G_{m,\ge r}:=\prod_{d_m(x)\ge r}\mathfrak S_{d_m(x)},
  \qquad
  N_m(r):=\#\{x:d_m(x)\ge r\}.
  \]
\end{enumerate}
若此猜想成立，则 rigid-window 理论将分裂为两层：其一是由 seed 完全支配的无标号半单热力学层；其二是必须依靠额外几何证书才能恢复的块定位层。window-\(6\) 的上述定理链给出了这一原则的完整锚点。
\end{conjecture}
